\section{Methodology}
\label{sec:method}

\subsection{Subjects and Protocol}
The study corpus consists of 63 Java functions drawn from a local transformation of the HumanEval benchmark \cite{yuan2024junit}. Each original Python function from HumanEval was re-implemented as a standalone Java class containing a correct implementation (the SUT) and a manually introduced buggy variant; functions with degenerate or trivially simple branch structure were excluded during corpus assembly. The correct implementation is the test target throughout this study. Prior work using the same benchmark in the Python setting, such as Yuan et al.\ \cite{yuan2024junit}, provides a useful point of reference, though the Java compilation and assertion requirements make direct score comparison between languages inappropriate.

For each target, the generation script supplied its complete source code to \texttt{gpt-4o-mini-2024-07-18} with a fixed zero-shot instruction. The prompt required a single executable JUnit~4 class in package \texttt{humaneval.correct}, correct assertions, and coverage of normal, boundary, and error cases. Temperature was fixed to zero. No output-token cap was set in the API call.

The run retained a timestamped API log and a CSV record for each target. The initial call was retried only for retryable service failures. One subsequent repair invocation was permitted for syntactic or assertion issues. Generated tests were then compiled and executed with Maven. Tests that remained invalid were archived rather than deleted, and their execution status was retained in the result set.

The comparator consists of existing EvoSuite archives generated under 1-, 3-, and 5-minute budgets. No EvoSuite generation was performed for this study. Each archive contains one suite per SUT and was measured against the same target set with the same Maven, JaCoCo, and PIT configuration. Temurin JDK~8 was used for this measurement because EvoSuite~1.0.6 requires \texttt{tools.jar}.

\subsection{Evidence Retention and Postprocessing}
The pipeline records a row per SUT for generation, repair, execution status, and class-level metrics. The primary files are the full API-output CSV, repair-output CSV, compile-status CSV, GPT class-metrics CSV, three EvoSuite class-metrics CSVs, a measurement manifest, and the derived summary CSV. The notebook reads these retained records rather than manually copying values into a report. A validation script checks the number of distinct SUTs, status partition, coverage numerators, mutation numerators, comparator pass counts, RQ3 rows, notebook outputs, and figure resolution.

Postprocessing does not repair the scientific record silently. Generated test classes are moved into timestamped archive or invalid-test locations after measurement rather than deleted. The active test directory is cleared only after preservation. This matters because a later reader must be able to distinguish a pipeline configuration change from a changed test suite. The paper reports the initial API run and the permitted repair pass separately, so the final execution gate is not confused with a single unaided model response.

\subsection{Measures and Decision Rules}
For every SUT, JaCoCo \cite{Lu_2026} produced branch-covered and branch-total counts; branch coverage is the covered-to-total ratio. PIT \cite{Bouafif_2025} produced killed-mutant and eligible-mutant counts; mutation score is their ratio. A suite was considered executable only if its test class compiled and passed all test cases without modification. The report retains zero or unavailable class-level outcomes for non-executable suites, while the paired RQ3 analysis includes only SUTs that passed in both tools.

RQ1 tests whether the per-SUT GPT branch score exceeds 30.22\%. RQ2 tests the mutation score against an empirical 4.00\% floor and a 40.21\% target. RQ4 tests whether more than half of the 63 SUTs simultaneously pass, reach branch coverage of at least 30.22\%, and reach mutation score of at least 4.00\%. RQ5 describes the frequency of technical failure patterns without an inferential test. RQ1 and RQ2 use one-sided Wilcoxon signed-rank tests; RQ4 uses an exact binomial test. The significance level is $\alpha=0.05$.

For RQ3, GPT and EvoSuite scores are paired by SUT for branch coverage and mutation score at each budget. The test is a two-sided paired Wilcoxon signed-rank test. Three technical mutation exclusions are omitted from mutation-score tests. The six p-values are Holm-adjusted. We report the raw and adjusted p-values, the paired and non-zero-ranked sample sizes, percentage-point difference, and rank-biserial effect size, following standard evaluation guidelines \cite{evaluating2024}. A comparison is considered statistically significant only when the adjusted p-value is below 0.05.

The report preserves two denominators for RQ3. The paired denominator counts SUTs that executed in both tools. The ranked denominator counts non-zero differences used by the signed-rank statistic. When many pairs have identical scores, the latter can be much smaller. Reporting both prevents a small ranked sample from being mistaken for 63 independent observations.

\subsection{Generation Prompt (Verbatim Template)}
\label{sec:prompt}
The following is the exact template passed as the user message. The two placeholders were instantiated separately for each SUT; \texttt{\{source\_code\}} was replaced by the complete source of that target. No example test, previous model response, or few-shot demonstration was included.

\begin{quote}\footnotesize
You are an expert Java developer and software tester. Your task is to write a comprehensive JUnit 4 test suite for the following Java class. Strictly adhere to the following requirements:\\
1. Generate test cases using JUnit 4 (use org.junit.Test, org.junit.Assert).\\
2. Do not use JUnit 5 or other test frameworks.\\
3. Test all logical paths, edge cases, boundary values, and potential error conditions.\\
4. Ensure all assertions are correct and correspond exactly to the expected behavior of the correct code.\\
5. The test class must be named \texttt{\{class\_name\}\_GPTTest}, corresponding to the target class \texttt{\{class\_name\}}.\\
6. The test class must be in package \texttt{humaneval.correct}.\\
7. You should import any required packages (like java.util.*).\\
8. Provide only the executable Java test class code. Do not include any markdown explanations, text wrapping, or extra commentary.\\
Source Code: \texttt{\{source\_code\}}
\end{quote}

\subsection{Reproducibility and Protocol Amendment}
The analysis notebook, derived CSVs, figures, API usage records, execution statuses, archive manifest, and validation report are stored in the replication package. The original proposal motivated a future comparison with student-written tests. Comparable per-SUT student measurements were unavailable, so that comparison is deferred. The present paper consequently makes no numerical claim about student performance and never treats EvoSuite as a student proxy.
